\begin{tikzpicture}[font=\small,line cap=round,line join=round]\draw[otMuted,dashed,line width=.9pt](0,0)--(2.2,0);\draw[otMuted,dashed,line width=.9pt](0,0)--(45:2.2);\node[below,font=\tiny]at(1.1,-.1){calque};\draw[geoTrait](3.4,0)--++(0:2.0);\draw[geoTrait](3.4,0)--++(30:2.0);\node[below,font=\small]at(4.3,-.25){$ D $};\draw[geoTrait](6.4,0)--++(0:2.0);\draw[geoTrait](6.4,0)--++(45:2.0);\node[below,font=\small]at(7.300000000000001,-.25){$ E $};\draw[geoTrait](9.4,0)--++(0:2.0);\draw[geoTrait](9.4,0)--++(70:2.0);\node[below,font=\small]at(10.3,-.25){$ F $};\end{tikzpicture}\par\vspace{1.5mm}{\otRegular\fontsize{8}{10}\selectfont\color{otMuted}Le calque en tirets se superpose à l’un des trois angles : même sommet, mêmes côtés.\par}
